\documentclass[10pt, aspectratio = 169]{beamer} %

\usetheme[progressbar=frametitle,numbering=fraction]{metropolis}
\usepackage{appendixnumberbeamer}

\usepackage{booktabs}
\usepackage[scale=2]{ccicons}
\usefonttheme[onlymath]{serif}
%\usepackage[sfdefault]{FiraSans}  % Metropolis default, very similar to Open Sans
\usepackage[default]{sourcesanspro}

\usepackage{pgfplots}
\usepgfplotslibrary{dateplot}
\usepackage{tikz}
\usepackage{graphicx}

\usepackage{subcaption}
\usepackage{multicol}

\usepackage{xspace}
\newcommand{\themename}{\textbf{\textsc{metropolis}}\xspace}

\usepackage{empheq}
\usepackage[many]{tcolorbox}

\usepackage{MnSymbol,wasysym}
\usepackage{hanging}
\setbeamertemplate{footnote}{\hangpara{2em}{1}\makebox[2em][l]{\insertfootnotemark}\footnotesize\insertfootnotetext\par}
\usetikzlibrary{fadings,positioning,calc,shadows}
\usepackage{setspace}
\usepackage[en-US]{datetime2}
\usepackage{picture}
\def\checkmark{\tikz\fill[scale=0.4](0,.35) -- (.25,0) -- (1,.7) -- (.25,.15) -- cycle;}
\newcommand{\Var}{\mathrm{Var}}
\newcommand{\Cov}{\mathrm{Cov}}
\newcommand{\yhat}{\widehat{y}}
\newcommand{\by}{\bm y}
\newcommand{\bZ}{(\bm X \bm C)}
\newcommand{\E}{\mathbb{E}}
\newcommand{\M}{\mathcal{M}}
\newcommand{\La}{\mathcal{L}}
\newcommand{\norm}[1]{\left\lVert#1\right\rVert}
\newcommand{\plim}{\xrightarrow[]{p}}
\newcommand{\dlim}{\xrightarrow[]{d}}
\newcommand{\infsum}{\sum_{t=0}^\infty}
\newcommand{\infint}{\int_{0}^\infty}
\newcommand{\intone}{\int_{0}^1}
\usetikzlibrary{arrows.meta,positioning}
\usepackage[capposition=top]{floatrow}
\usepackage{subcaption}
\usepackage{booktabs}
\usepackage{array}

\usepackage{jambro-annotations}
\definecolor{carandache}{RGB}{144,26,30}

%%%%%%%%%%%%%%%%%%%%%%%%%%%%%%%%%%%%%%%%%%%%%%%%%%%%%%%%%%%%%%
%% Color pallette
\definecolor{cph_red}{RGB}{144,26,30}
\definecolor{cph_grey}{RGB}{102,102,102}
\setbeamercolor{palette primary}{fg=cph_red, bg= white} %{fg=white, bg= cph_red}

\setbeamercolor{button}{bg=cph_red}
\setbeamercolor{title}{fg=black}
\setbeamercolor{section title}{fg=cph_red!100}
\setbeamercolor{progress bar}{fg=cph_grey, bg= white}
\setbeamercolor{itemize item}{fg=cph_red,bg=white}

\setbeamercolor{alerted text}{fg=cph_red, bg= white}

% Change annotation color to cph red
\definecolor{carandache}{RGB}{0,0,128}   % define your own

\makeatletter
\setlength{\metropolis@titleseparator@linewidth}{2pt}
\setlength{\metropolis@progressonsectionpage@linewidth}{2pt}
\setlength{\metropolis@progressinheadfoot@linewidth}{2pt}
\makeatother

%%%%%%%%%%%%%%%%%%%%%%%%%%%%%%%%%%%%%%%%%%%%%%%%%%%%%%%%%%%%%

\tcbset{highlight math style={enhanced,
  colframe=red!60!black,colback=yellow!50!white,arc=4pt,boxrule=1pt,
  }}


\tikzset{
    bluenode/.style={
        draw,
        black,
        top color=orange!50!white,
        bottom color=orange!90!white,
        rounded corners,drop shadow
    },
    graynode/.style={
        draw,
        top color=gray!50!white!20,
        bottom color=gray!50!black!20,drop shadow
    }}
\title{The Curious Incidence of Monetary Policy Across the Income
Distribution}

\date{\today}
\author{Tobias Broer \ \ \   John Kramer \ \ \  Kurt Mitman}
\DTMlangsetup{showdayofmonth=false}

\makeatletter
\def\th@mystyle{%
    \normalfont % body font
    \setbeamercolor{block title example}{bg=orange,fg=white}
    \setbeamercolor{block body example}{bg=orange!20,fg=black}
    \def\inserttheoremblockenv{exampleblock}
  }
\makeatother
\theoremstyle{mystyle}
\newtheorem*{conjecture}{Conjecture}

% From stata regression table
\def\sym#1{\ifmmode^{#1}\else\(^{#1}\)\fi}

%\institute{\today}
% \titlegraphic{\hfill\includegraphics[height=1.5cm]{logo.pdf}}

\begin{document}
%%%%%%%%%%%%%%%%%%%%%%%%%%%%%%%%%%%%%%%%%%%%%%%%%%%%%%%%%%%%%%%%%%%%%%%%
\setbeamercolor{progress bar}{fg=cph_red, bg= white}
%%%%%%%%%%%%%%%%%%%%%%%%%%%%%%%%%%%%%%%%%%%%%%%%%%%%%%%%%%%%%%%%%%%%%%%%
\begin{frame}[noframenumbering, plain]
\maketitle
\begin{tikzpicture}[overlay, remember picture]
\node[above left=0.3cm and 0.6cm of current page.south east]{\includegraphics[width=5cm,trim = 31cm 10cm 1cm 15cm,clip]{Presentations/KU_RGB.pdf}};
\end{tikzpicture}
\end{frame}
%%%%%%%%%%%%%%%%%%%%%%%%%%%%%%%%%%%%%%%%%%%%%%%%%%%%%%%%%%%%%%%%%%%%%%%%
\setbeamercolor{progress bar}{fg=cph_grey, bg= white}



%%%%%%%%%%%%%%%%%%%%%%%%%%%%%%%%%%%%%%%%%%%%%%%%%%%%%%%%%%%%%%%%%%%%%%%%
\begin{frame}{Motivation}
\onslide<1->{\alert{Understanding the effects of monetary policy}
\begin{itemize}
    \item Recent aggressive monetary tightening raises distributional concerns
    \item Limited evidence on who bears the cost of monetary policy interventions
    \item Theory suggests heterogeneity may amplify or dampen aggregate effects
\end{itemize}}
\onslide<2>{{\centering \large \alert{Key questions:}\par}
\begin{itemize}
    \item How does monetary policy affect individuals across the income distribution? 
    \item Does this heterogeneity matter for transmission of MP to  aggregate demand? 
\end{itemize}
$\implies$ empirical \& theoretical analysis}
\end{frame}
%%%%%%%%%%%%%%%%%%%%%%%%%%%%%%%%%%%%%%%%%%%%%%%%%%%%%%%%%%%%%%%%%%%%%%%%


%%%%%%%%%%%%%%%%%%%%%%%%%%%%%%%%%%%%%%%%%%%%%%%%%%%%%%%%%%%%%%%%%%%%%%%%
\begin{frame}{Empirical Analysis}

\alert{Approach}
\begin{itemize}
    \item Using German administrative data, we estimate responses to monetary policy
\begin{itemize}
     \item for earnings growth
     \item for labor market transitions probabilities
\end{itemize}
across the permanent income distribution
\end{itemize}

\onslide<2->{\alert{Findings}
\begin{itemize}
    \item Monetary policy impacts earnings growth at the bottom twice as much as at the top
    \item Heterogeneity is driven by extensive margin transitions
    \item Employed workers show uniform earnings responses
    \item Effects are long-lasting for those who become unemployed
\end{itemize}}


\end{frame}
%%%%%%%%%%%%%%%%%%%%%%%%%%%%%%%%%%%%%%%%%%%%%%%%%%%%%%%%%%%%%%%%%%%%%%%%



%%%%%%%%%%%%%%%%%%%%%%%%%%%%%%%%%%%%%%%%%%%%%%%%%%%%%%%%%%%%%%%%%%%%%%%%
\begin{frame}{Theoretical Framework}

\alert{Model}
\begin{itemize}
\item HANK model to study transmission of monetary policy
\item Permanent type heterogeneity in permanent income
\item Heterogeneous and counter-cyclical unemployment risk 
\end{itemize}

\onslide<2>{\alert{Findings}
\begin{itemize}
    \item Heterogeneous and counter-cyclical \lapis[cph_red]{unemployment} risk amplifies effects of monetary policy by 20\% relative to pure \lapis[cph_red]{income} risk
    \begin{itemize}
        \item Precautionary savings from unemployment risk drive half of the amplification
        \item Rest is driven by realized income differences
    \end{itemize}
\end{itemize}
\textbf{Bottom line:} Labor income heterogeneity is crucial for understanding monetary policy effectiveness }
\end{frame}
%%%%%%%%%%%%%%%%%%%%%%%%%%%%%%%%%%%%%%%%%%%%%%%%%%%%%%%%%%%%%%%%%%%%%%%%


%%%%%%%%%%%%%%%%%%%%%%%%%%%%%%%%%%%%%%%%%%%%%%%%%%%%%%%%%%%%%%%%%%%%%%%%
\begin{frame}{Literature and contribution}


\alert{Empirical evidence on heterogeneous effects of monetary policy}
\begin{itemize}
    \item \footnotesize{Coibion et al (2017), Guvenen et al (2017), D’Acunto et al (2019a,b), Patterson (2018), De Giorgi \& Gambetti (2017), Alves et al (2019), Almgren et al (2019), Holm et al (2021), \lapis[cph_red]{Amberg et al (2022)}, \lapis[cph_red]{Andersen et al (2023)}, Hubert \& Savignac (2024),...}
\end{itemize}
\alert{What we add:} Heterogeneous earnings responses driven by \textbf{unemployment}
\vspace{0.5cm}

\alert{(Quantitative) Theory:}
\begin{itemize}
    \item \footnotesize{Gornemann et al (2012), Auclert (2019), Kaplan et al (2018), \lapis[cph_red]{Werning (2015)}, Broer et al (2019), Hagedorn et al (2019a,b), {Bilbiie (2018)}, \lapis[cph_red]{Ravn \& Sterk (2017, 2021)}, {Acharya and Dogra (2020)}, Auclert et al (2021), Bilbiie et al (2022), \lapis[cph_red]{Graves (2025)},...}
\end{itemize}
\alert{What we add:} \textbf{Heterogeneous}, countercyclical \textbf{unemployment risk} amplifies output responses
\end{frame}
%%%%%%%%%%%%%%%%%%%%%%%%%%%%%%%%%%%%%%%%%%%%%%%%%%%%%%%%%%%%%%%%%%%%%%%%


%%%%%%%%%%%%%%%%%%%%%%%%%%%%%%%%%%%%%%%%%%%%%%%%%%%%%%%%%%%%%%%%%%%%%%%%
\section{Data}
%%%%%%%%%%%%%%%%%%%%%%%%%%%%%%%%%%%%%%%%%%%%%%%%%%%%%%%%%%%%%%%%%%%%%%%%
\begin{frame}{German administrative micro-data}
Social security data from German federal employment agency (SIAB sample)
\begin{itemize}
    \item 2\% of all individuals in social-insurance system (excl public employees \& self-employed)
    \item Employment history \& ``Daily'' employment status (agg to \lapis[cph_red]{monthly})
    \item Main variables of interest:\\
    \vspace{2pt}
    - \lapis[cph_red]{Employment} spells (start and end date)\\
    \vspace{2pt}
    - Within-spell average pre-tax labor \lapis[cph_red]{earnings} (impute top coded$\sim 6$ \%)
    \item Restrict to closely attached: employed \& unemployed 
    \item Sample period: 2000-2013 (excl zero lower bound)
\end{itemize}
\alert{$\to$ Allows for clean decomposition of earnings movements}
\end{frame}
%%%%%%%%%%%%%%%%%%%%%%%%%%%%%%%%%%%%%%%%%%%%%%%%%%%%%%%%%%%%%%%%%%%%%%%%


%%%%%%%%%%%%%%%%%%%%%%%%%%%%%%%%%%%%%%%%%%%%%%%%%%%%%%%%%%%%%%%%%%%%%%%%
\begin{frame}{Heterogeneous effects along the income distribution}\label{hetintro}
\alert{Income distribution}  \hyperlink{descr}{\beamergotobutton{Descriptive statistics}}
\begin{itemize}
    \item Sort individuals in \textbf{5-year-average earnings} quantiles $q$ conditional on age and gender
    \item Similar results for other measures of permanent income
\end{itemize}

\onslide<2>{\alert{Effects of monetary policy} \hyperlink{aggs0}{\beamergotobutton{Aggregates}}
\begin{align*}
    \ln \left(\overline{earn}_{t+k} \right) - \ln \left(\overline{earn}_{t-1} \right) = \alpha^q_{k} + \beta^q_{k} \Delta i_{t} +\theta^q_{k}  X_{t}+\epsilon^q_{t+k}\quad \text{for} \quad k = 12
\end{align*}
\vspace{-0.5cm}
\begin{itemize}
    \item Instrument policy rate w/ OIS rate changes around announcement (Almgren et al. 2022)
    \item Separate IV estimates for individuals in quantile $q$ in $t-1$
    \item Dependent variables $x^q_{t+k}$:
    \begin{itemize}
        \item Change in \textbf{quantile earnings} from $t$ to $t+k$
    \end{itemize}
    \item Scale $\Delta i_{t}$ to cause 1pp rise in average earnings $\overline{earn}_t$
\end{itemize}}
\end{frame}
%%%%%%%%%%%%%%%%%%%%%%%%%%%%%%%%%%%%%%%%%%%%%%%%%%%%%%%%%%%%%%%%%%%%%%%%

\section{Empirical results}

%%%%%%%%%%%%%%%%%%%%%%%%%%%%%%%%%%%%%%%%%%%%%%%%%%%%%%%%%%%%%%%%%%%%%%%%
\begin{frame}{Monetary policy incidence on earnings growth across the distribution}
\begin{columns}[c]
\begin{column}{0.6\textwidth}
    \centering
    \includegraphics[width=\linewidth]{Figures/earn_n_tot.pdf} 
\end{column}
\hspace{-10mm}
\begin{column}{0.4\textwidth}
    \begin{itemize}
        \item \textbf{Earnings growth most affected at the bottom} 
        
        \footnotesize $\sim$ Patterson (2023), Amberg et al (2022)
        \vspace{.5cm}
        \item \normalsize Uptick at the top likely due to imputation
        \vspace{0.5cm}
        \item \normalsize Monetary policy impact similar to unconditional 
        
         \footnotesize $\sim$ Guvenen et al (2017)
    \end{itemize}
\end{column}
\end{columns}
\end{frame}
%%%%%%%%%%%%%%%%%%%%%%%%%%%%%%%%%%%%%%%%%%%%%%%%%%%%%%%%%%%%%%%%%%%%%%%%

%%%%%%%%%%%%%%%%%%%%%%%%%%%%%%%%%%%%%%%%%%%%%%%%%%%%%%%%%%%%%%%%%%%%%%%%
\begin{frame}{Decomposition into earnings and employment}\label{decomposition}
\begin{columns}[c]
\begin{column}{0.6\textwidth}
    \centering
    \includegraphics[width=\linewidth]{Figures/earn_n_e2e_dash.pdf} 
\end{column}
\hspace{-10mm}
\begin{column}{0.4\textwidth}
    \begin{align*}
    &\ln \left(\overline{earn}_{t} \right) = \ln \left(s^{emp}_t \cdot\overline{earn}^e_{t} \right) = \\
    &\ln\left( s^{emp}_t \right) + \ln \left(\overline{earn}^e_{t} \right)
\end{align*}
    \begin{itemize}
        \item Earnings of the employed homogeneously procyclical  
        \hyperlink{semp}{\beamergotobutton{Complement}} \hyperlink{UI}{\beamergotobutton{with UI}} 
    \end{itemize}
\end{column}
\end{columns}

\vspace{5mm}
\centering
\alert{Unemployment is driving higher incidence of monetary policy on the poor}

\end{frame}
%%%%%%%%%%%%%%%%%%%%%%%%%%%%%%%%%%%%%%%%%%%%%%%%%%%%%%%%%%%%%%%%%%%%%%%%



%%%%%%%%%%%%%%%%%%%%%%%%%%%%%%%%%%%%%%%%%%%%%%%%%%%%%%%%%%%%%%%%%%%%%%%%
\begin{frame}{How does monetary policy affect (un)employment?}
\alert{Compute transition probabilities}
\begin{itemize}
    \item Define transition probability as share of individuals in quantile $q$ in state $s_{t-1}\in\{U,E\}$ who move to state $s_{t+k}\in\{U,E \}$
    \item[$\to$] Times series of transition probabilities by quantile
\end{itemize}
\vspace{0.5cm}
\alert{Effect of monetary policy on labor market transitions}
\begin{align*}
    TR^{q,s_1, s_2}_{t+k} = \alpha^q_{k} + \beta^q_{k} \Delta i_{t} +\theta^q_{k}  X_{t}+\epsilon^q_{t+k}\quad \text{for} \quad k = 12
\end{align*}
\end{frame}
%%%%%%%%%%%%%%%%%%%%%%%%%%%%%%%%%%%%%%%%%%%%%%%%%%%%%%%%%%%%%%%%%%%%%%%%



%%%%%%%%%%%%%%%%%%%%%%%%%%%%%%%%%%%%%%%%%%%%%%%%%%%%%%%%%%%%%%%%%%%%%%%%
\begin{frame}{Incidence of monetary policy on transition probabilities}\label{transprob}
\begin{columns}[c]
\begin{column}{0.49\textwidth}
    \centering
    \alert{Separations} --- $s_{t-1}=E; s_{t+12} =U$
    
    \vspace{3pt}
    \includegraphics[width=\linewidth]{Figures/trans_n_e2u.pdf}
    \begin{itemize}
        \item Separations fall for bottom, top unaffected
    \end{itemize}
\end{column}
\begin{column}{0.49\textwidth}
    \centering
    \alert{Job-finding} --- $s_{t-1}=U; s_{t+12} =E$
    
    \vspace{3pt}
    \includegraphics[width=\linewidth]{Figures/trans_n_u2e.pdf}
    \begin{itemize}
        \item Impact on job-finding homogeneous \hyperlink{switching}{\beamergotobutton{Job switching}}

    \end{itemize}
\end{column}
\end{columns}
\end{frame}
%%%%%%%%%%%%%%%%%%%%%%%%%%%%%%%%%%%%%%%%%%%%%%%%%%%%%%%%%%%%%%%%%%%%%%%%


%%%%%%%%%%%%%%%%%%%%%%%%%%%%%%%%%%%%%%%%%%%%%%%%%%%%%%%%%%%%%%%%%%%%%%%%
\begin{frame}{Taking stock}\label{stock}
\alert{Empirical results} \ \ \ \hyperlink{LRempl}{\beamergotobutton{Long-run employment}} \hyperlink{LRearn}{\beamergotobutton{Long-run earnings}} \hyperlink{gini}{\beamergotobutton{Inequality}}
\begin{itemize}
    \item Earnings growth at the bottom of the income distribution is most affected by monetary policy
    \item The main driver behind the heterogeneity are \textbf{extensive margin transitions} 
    \item \textbf{Separation probabilities are very counter-cyclical} at the bottom of the distribution
\end{itemize}
\alert{Two implications}
\vspace{2pt}
\begin{enumerate}
    \item Less productive workers have more cyclical earnings growth rates
\item Due to both higher \underline{and} more cyclical unemployment risk
\end{enumerate}
\vspace{5pt}
\textbf{Rest of the talk}\\
\vspace{2pt}
\hspace{5pt} Show how this matters for the transmission of monetary policy
\end{frame}
%%%%%%%%%%%%%%%%%%%%%%%%%%%%%%%%%%%%%%%%%%%%%%%%%%%%%%%%%%%%%%%%%%%%%%%%

\section{Model}
%%%%%%%%%%%%%%%%%%%%%%%%%%%%%%%%%%%%%%%%%%%%%%%%%%%%%%%%%%%%%%%%%%%%%%%%
\begin{frame}{Model setup}\label{modelsetup}
\alert{Heterogeneity by permanent income}
\begin{itemize}
    \item Introduce ex-ante, permanent heterogeneity 
    \item Types $g$ differ by productivity, discount factor and unemployment risk 
\end{itemize}

\onslide<2->{\alert{Earnings risk}
\begin{itemize}
    \item Common, persistent income process when employed ($x_{i,t}$)
    \item Exogenous, type specific unemployment probability
\end{itemize}
}
\onslide<3->{
\alert{Heterogeneous, countercyclical unemployment risk}
\begin{itemize}
    \item Exogenous feedback from aggregates to unemployment risk
\end{itemize}

Otherwise standard quarterly HANK model \hyperlink{firmdetails}{\beamergotobutton{Details}}
}
\end{frame}
%%%%%%%%%%%%%%%%%%%%%%%%%%%%%%%%%%%%%%%%%%%%%%%%%%%%%%%%%%%%%%%%%%%%%%%%

%%%%%%%%%%%%%%%%%%%%%%%%%%%%%%%%%%%%%%%%%%%%%%%%%%%%%%%%%%%%%%%%%%%%%%%%
\begin{frame}{Income risk for households}
\alert{Households utility and budget}
\begin{align*}
\max \quad \mathbb{E}_0 \sum_{t=0}^{\infty} \beta_g^t
\left(
\frac{c_{i,t}^{1-\gamma}}{1-\gamma}\right) \quad \text{ s.t. } \quad 
c_{i,t} + a_{i,t+1} =\marker{return}{$(1+r^p_t)$}a_{i,t}+(1-\tau)y_{i,t}+\marker{transfer}{$T_t$}; \quad \quad a_{i,t} \geq 0
\end{align*}
\onslide<3->{\alert{Labor income}
\vspace{-0.3cm}
\begin{align*}
    y^e_{i,t} = w_t h_{g,t} z_g x_{i,t}
\end{align*}
\vspace{-0.7cm}
\begin{itemize}
    \item Varies because of the real wage $w_t$, type specific hours $h_{g,t}$, or individual productivity $x_{i,t}$ 
    \item Replacement rate in unemployment is 60\%, as in the data
\end{itemize}}
\onslide<4->{
\alert{Unemployment}
\vspace{-0.3cm}
\begin{align*}
    u_{g,t} = \delta_{g,t}(1-u_{g,t-1})+(1-f_{g,t})u_{g,t-1}
\end{align*}
\vspace{-0.7cm}
\begin{itemize}
    \item Separation rate $\delta_{g,t}$ and job-finding rate $f_{g,t}$ are cyclical, as in the data
\end{itemize}
}
% Add annotations
\onslide<2>{\lapisnote[from=north,color=cph_red, yshift=0.6cm, xshift=0.5cm]{return}{Return on mutual fund}
\lapisnote[from=north east,color=cph_red, xshift=0.6cm, yshift=0.3cm]{transfer}{Transfers}}

\end{frame}
%%%%%%%%%%%%%%%%%%%%%%%%%%%%%%%%%%%%%%%%%%%%%%%%%%%%%%%%%%%%%%%%%%%%%%%%



%%%%%%%%%%%%%%%%%%%%%%%%%%%%%%%%%%%%%%%%%%%%%%%%%%%%%%%%%%%%%%%%%%%%%%%%
\begin{frame}{Counter-cyclical risk }
Baseline: \alert{``Endogenous'', type-specific separation and job-finding rates } 
\begin{align*}
  \delta^g(w) &= d^g_0 + d^g_1 \left(\log(w)-\log(w_{ss})\right) \\
  f^g(w) &= f^g_0 + f^g_1 \left(\log(w)-\log(w_{ss})\right) 
\end{align*}
\vspace{-.5cm}
\begin{itemize}
    \item Exogenous wage Phillips curve: 
    \begin{itemize}
        \item real wage drives labor-input fluctuations
    \end{itemize}
    \item Simple setup makes counterfactual implementation easy 
    \item Steady state deviations make things more stable
\end{itemize}

\end{frame}
%%%%%%%%%%%%%%%%%%%%%%%%%%%%%%%%%%%%%%%%%%%%%%%%%%%%%%%%%%%%%%%%%%%%%%%%

%%%%%%%%%%%%%%%%%%%%%%%%%%%%%%%%%%%%%%%%%%%%%%%%%%%%%%%%%%%%%%%%%%%%%%%%
\begin{frame}{Calibration}\label{calibration}

\alert{Discount factors} 
\begin{itemize}
    \item Target the same MPC distribution across all counterfactuals \hyperlink{table3}{\beamergotobutton{Internally calibrated parameters}}
    \item MPCs by permanent income type: [0.30, 0.25, 0.20] (Fagereng et al. 2021; Boehm et al. 2025)
\end{itemize}

\alert{Earnings risk}
\begin{itemize}
    \item Persistent risk of the employed following Krueger et al. (2016) 
    \item Transition probability process and productivities from the data \hyperlink{table2}{\beamergotobutton{Transition parameters}}
\end{itemize}

\alert{Other parameters}
\begin{itemize}
    \item Standard HANK calibration \hyperlink{table1}{\beamergotobutton{Externally calibrated parameters}}
\end{itemize}

\end{frame}
%%%%%%%%%%%%%%%%%%%%%%%%%%%%%%%%%%%%%%%%%%%%%%%%%%%%%%%%%%%%%%%%%%%%%%%%

%%%%%%%%%%%%%%%%%%%%%%%%%%%%%%%%%%%%%%%%%%%%%%%%%%%%%%%%%%%%%%%%%%%%%%%%
\begin{frame}{Baseline responses to contractionary monetary policy}\label{Cirf}
\begin{columns}[c]
\begin{column}{0.6\textwidth}
    \centering
    \includegraphics[trim=0cm 6cm 0cm 6.5cm, clip, width=\linewidth]{modelresults/single_CTlev_bgy397_tp0_ui6_sr0.pdf}
\end{column}
\hspace{-10mm}
\begin{column}{0.4\textwidth}
\begin{itemize}
    \item 25BP monetary policy shock reduces consumption by 1.4\%
\vspace{0.5cm}
    \item Unemployment rises by about 1 percentage point
\vspace{0.5cm}
    \item The real wage falls  \hyperlink{baseirf1}{\beamergotobutton{Other irfs}}
\end{itemize}
\end{column}
\end{columns}

\centering
\alert{Question:} How much is driven by unemployment risk?

\end{frame}
%%%%%%%%%%%%%%%%%%%%%%%%%%%%%%%%%%%%%%%%%%%%%%%%%%%%%%%%%%%%%%%%%%%%%%%%

%%%%%%%%%%%%%%%%%%%%%%%%%%%%%%%%%%%%%%%%%%%%%%%%%%%%%%%%%%%%%%%%%%%%%%%%
\begin{frame}{Counterfactual economies}
\alert{Four counterfactuals}
\begin{itemize}
    \item Unemployment transitions or hours adjust in response to shocks
    \item Incidence is heterogeneous or homogeneous in steady state and along transition
\end{itemize}
\begin{center}
\begin{tabular}{>{\centering\arraybackslash}m{2.2cm}
                >{\centering\arraybackslash}m{3.2cm}
                >{\centering\arraybackslash}m{3.2cm}}
    & \multicolumn{2}{c}{\textbf{Incidence}} \\
    \cmidrule(lr){2-3}
    \textbf{Margin} 
        & \textbf{\tikz[remember picture] \node (hetHead) {Heterogeneous};} 
        & \textbf{Homogeneous} \\
    \midrule
    \textbf{Unemployment} 
        & \tikz[remember picture] \node (hetU) {\alert{Heterogeneous U}}; 
        & Homogeneous U \\
    \midrule
    \textbf{Hours} 
        & \tikz[remember picture] \node (hetH) {Heterogeneous H}; 
        & Homogeneous H \\
    \bottomrule
\end{tabular}
\end{center}
\onslide<2>{\begin{tikzpicture}[remember picture, overlay]
\node[pencil,draw=cph_red, thick,
      fit=(hetHead) (hetU) (hetH),
      inner sep=6pt] { };
\end{tikzpicture}}

\alert{Keep things comparable}
\begin{itemize}
    \item Relationship between real wage movements and total labor input constant 
    \item For homogeneous transition probabilities, re-compute elasticities
\end{itemize}

\end{frame}
%%%%%%%%%%%%%%%%%%%%%%%%%%%%%%%%%%%%%%%%%%%%%%%%%%%%%%%%%%%%%%%%%%%%%%%%



%%%%%%%%%%%%%%%%%%%%%%%%%%%%%%%%%%%%%%%%%%%%%%%%%%%%%%%%%%%%%%%%%%%%%%%%
\begin{frame}{Hours vs unemployment and the importance of risk}
\begin{columns}[c]
\begin{column}{0.6\textwidth}
    \centering
    \includegraphics[trim=0cm 6cm 0cm 6.5cm, clip, width=\linewidth]{modelresults/Consumption_Comparison.pdf}
\end{column}
\hspace{-10mm}
\begin{column}{0.4\textwidth}
\begin{itemize}
    \item Heterogeneous risk loaded on hours attenuates response

    $\to$ \lapis[cph_red]{unemployment matters}
\vspace{0.5cm}
    
    \onslide<2>{\item Does \textbf{risk} matter? 
    
    \item[$\to$] Economy with expected unemployment but hours movements:
    $C_t^{risk} = \int \Omega^{Het H} C_{i}^{Het U} \ di$
\vspace{0.5cm}
    \item Unemployment \lapis[cph_red]{risk} drives half of the difference}
\end{itemize}
\end{column}
\end{columns}

\centering
\onslide<2>{\alert{Unemployment risk amplifies the effects of monetary policy}}

\end{frame}
%%%%%%%%%%%%%%%%%%%%%%%%%%%%%%%%%%%%%%%%%%%%%%%%%%%%%%%%%%%%%%%%%%%%%%%%


%%%%%%%%%%%%%%%%%%%%%%%%%%%%%%%%%%%%%%%%%%%%%%%%%%%%%%%%%%%%%%%%%%%%%%%%
\begin{frame}{The importance of heterogeneous incidence}\label{hetinc0}
\begin{columns}[c]
\begin{column}{0.6\textwidth}
    \centering
    \includegraphics[trim=0cm 6cm 0cm 6.5cm, clip, width=\linewidth]{modelresults/aggs3_1_CTlev_bgy397_tp0_ui6_sr.pdf}
\end{column}
\hspace{-10mm}
\begin{column}{0.4\textwidth}
\begin{itemize}
    \item \lapis[cph_red]{Homogeneity attenuates} consumption response
    \vspace{0.3cm}

    \item[$\to$] loads more unemployment on highly productive workers, hence smaller $U$ movements 

    \vspace{0.3cm}
    \item[$\to$] smaller income fluctuations for poor, high MPC agents
    
    \vspace{0.3cm}
    \item[$\to$] Different steady state \hyperlink{hetinc}{\beamergotobutton{Decomposition}}
    
\end{itemize}
\end{column}
\end{columns}

\centering
\onslide<2>{\alert{Heterogeneity is important in driving consumption response}}

\end{frame}
%%%%%%%%%%%%%%%%%%%%%%%%%%%%%%%%%%%%%%%%%%%%%%%%%%%%%%%%%%%%%%%%%%%%%%%%




%%%%%%%%%%%%%%%%%%%%%%%%%%%%%%%%%%%%%%%%%%%%%%%%%%%%%%%%%%%%%%%%%%%%%%%%
\begin{frame}{Conclusion}
     \alert{Monetary policy has significantly larger effects on low-income workers' earnings}
    \begin{itemize}
        \item Heterogeneity driven primarily by unemployment risk at the bottom
        \item No heterogeneity in earnings response of the employed 
    \end{itemize}
    
    \vspace{0.3cm}
    
     \alert{Heterogeneous unemployment risk amplifies monetary policy transmission}
    \begin{itemize}
        \item Precautionary savings channel leads to large consumption drop
        \item Heterogeneity loads on high-MPC workers
        \item Consumption response increases by about one-third relative to homogeneous-risk models
    \end{itemize}
    
    \vspace{0.3cm}
    
     \textbf{Bottom line:} Labor income heterogeneity is crucial for understanding monetary policy effectiveness and its distributional consequences

\end{frame}
%%%%%%%%%%%%%%%%%%%%%%%%%%%%%%%%%%%%%%%%%%%%%%%%%%%%%%%%%%%%%%%%%%%%%%%%



\appendix

%%%%%%%%%%%%%%%%%%%%%%%%%%%%%%%%%%%%%%%%%%%%%%%%%%%%%%%%%%%%%%%%%%%%%%%%
\begin{frame}{Sample aggregates}\label{aggs0}

\alert{Aggregate effects} 
\vspace{-3mm}
\begin{align*}
    \Delta x_{t+k}= \alpha_{k} + \beta_{k}\Delta i_{t}     +\theta_{k}    X_{t}+\epsilon_{t+k}
\end{align*}



\begin{columns}[t]
\begin{column}{0.48\textwidth}
    \centering
    \footnotesize
    Aggregate earnings growth 
    
    $ \Delta x_{t+k}=\ln(\overline{earn}_{t+k})-\ln(\overline{earn}_{t-1})$ 
    \vspace{1mm}
    \includegraphics[width=.9\linewidth]{Figures/agg_n_agg_earn.pdf}
\end{column}
\begin{column}{0.48\textwidth}
    \centering
    \footnotesize
    Probability of staying employed
    
    $\Delta x_{t+k}=\overline{emp}_{t+k}|emp_{it-1}=1$ 
    \vspace{1mm}
    \includegraphics[width=.9\linewidth]{Figures/agg_n_employed_fa.pdf}
\end{column}
\end{columns}

\hyperlink{hetintro}{\beamergotobutton{Back}}

\end{frame}
%%%%%%%%%%%%%%%%%%%%%%%%%%%%%%%%%%%%%%%%%%%%%%%%%%%%%%%%%%%%%%%%%%%%%%%%

%%%%%%%%%%%%%%%%%%%%%%%%%%%%%%%%%%%%%%%%%%%%%%%%%%%%%%%%%%%%%%%%%%%%%%%%
\begin{frame}{Aggregate effects}\label{aggs}

\begin{columns}[t]
\begin{column}{0.48\textwidth}
    \centering
    \small
    Inflation
    
    \includegraphics[width=.7\linewidth]{Figures/agg_n_cpi.pdf}
    
    Unemployment rate
    
    \includegraphics[width=.7\linewidth]{Figures/agg_n_u_rate.pdf}
\end{column}
\begin{column}{0.48\textwidth}
    \centering
    \small
    Industrial production
    
    \includegraphics[width=.7\linewidth]{Figures/agg_n_IP.pdf}
    
    Real rate
    
    \includegraphics[width=.7\linewidth]{Figures/agg_n_r.pdf}
\end{column}
\end{columns}
\hfill \hyperlink{hetintro}{\beamergotobutton{Back}}

\end{frame}
%%%%%%%%%%%%%%%%%%%%%%%%%%%%%%%%%%%%%%%%%%%%%%%%%%%%%%%%%%%%%%%%%%%%%%%%



%%%%%%%%%%%%%%%%%%%%%%%%%%%%%%%%%%%%%%%%%%%%%%%%%%%%%%%%%%%%%%%%%%%%%%%%
\begin{frame}{Descriptive statistics by decile}\label{descr}
\centering
\resizebox{\textwidth}{!}{\input{Figures/graph2010fa_new}}

\hyperlink{hetintro}{\beamergotobutton{Back}}
\end{frame}
%%%%%%%%%%%%%%%%%%%%%%%%%%%%%%%%%%%%%%%%%%%%%%%%%%%%%%%%%%%%%%%%%%%%%%%%


%%%%%%%%%%%%%%%%%%%%%%%%%%%%%%%%%%%%%%%%%%%%%%%%%%%%%%%%%%%%%%%%%%%%%%%%
\begin{frame}{Externally calibrated parameters
}\label{table1}
\centering
\resizebox{0.7\textwidth}{!}{\input{tex/calibrationtable}}

\hyperlink{calibration}{\beamergotobutton{Back}}
\end{frame}
%%%%%%%%%%%%%%%%%%%%%%%%%%%%%%%%%%%%%%%%%%%%%%%%%%%%%%%%%%%%%%%%%%%%%%%%

%% Slide 1: Supply Side
%%%%%%%%%%%%%%%%%%%%%%%%%%%%%%%%%%%%%%%%%%%%%%%%%%%%%%%%%%%%%%%%%%%%%%%%
\begin{frame}{Model: Supply Side}\label{firmdetails}

\textbf{Final Good Producer}
\[
Y_t = \left( \int_0^1 Y_{j,t}^{\frac{\varepsilon-1}{\varepsilon}} dj \right)^{\frac{\varepsilon}{\varepsilon-1}}, \quad P_t = \left( \int_0^1 p_{j,t}^{1-\varepsilon} dj \right)^{\frac{1}{1-\varepsilon}}
\]

\textbf{Intermediate Goods Firms}
\begin{itemize}
    \item Linear production: $Y_{j,t} = H_{j,t}$, marginal cost $mc_{j,t} = w_t$
    \item Rotemberg pricing $\Rightarrow$ standard NKPC
    \item Fixed operating cost $F$ (zero steady-state profits)
\end{itemize}

\textbf{Mutual Fund}
\begin{itemize}
    \item Collects household savings, invests in government bonds
    \item Owns firms, distributes profits: $r^p_t = r_t + \pi_t / A_{t-1}$
\end{itemize}
\hyperlink{modelsetup}{\beamergotobutton{Back}}

\end{frame}
%%%%%%%%%%%%%%%%%%%%%%%%%%%%%%%%%%%%%%%%%%%%%%%%%%%%%%%%%%%%%%%%%%%%%%%%

%% Slide 2: Government & Equilibrium
%%%%%%%%%%%%%%%%%%%%%%%%%%%%%%%%%%%%%%%%%%%%%%%%%%%%%%%%%%%%%%%%%%%%%%%%
\begin{frame}{Model: Government \& Equilibrium}

\textbf{Fiscal Policy}
\[
B_t = (1+r_t)B_{t-1} + G_t + (1-\tau) UI_t + T_t - \tau W_t H_t
\]
Transfers $T_t$ adjust to balance the budget each period.

\textbf{Monetary Policy (Taylor Rule)}
\[
i_t = \bar{i} + \phi_\pi \pi_t + \varepsilon_t^{mp}
\]

\textbf{Market Clearing}
\begin{itemize}
    \item Goods: $Y_t = C_t + G_t - F$
    \item Assets: $A_t = B_t$
    \item Distribution: $\Omega_{t+1} = \mathcal{H}(\Omega_t)$
\end{itemize}
\hyperlink{modelsetup}{\beamergotobutton{Back}}

\end{frame}
%%%%%%%%%%%%%%%%%%%%%%%%%%%%%%%%%%%%%%%%%%%%%%%%%%%%%%%%%%%%%%%%%%%%%%%%


%%%%%%%%%%%%%%%%%%%%%%%%%%%%%%%%%%%%%%%%%%%%%%%%%%%%%%%%%%%%%%%%%%%%%%%%
\begin{frame}{Parameters governing unemployment risk 
}\label{table2}
\centering
\resizebox{0.7\textwidth}{!}{\input{tex/elasticity_table}}

\hyperlink{calibration}{\beamergotobutton{Back}}
\end{frame}
%%%%%%%%%%%%%%%%%%%%%%%%%%%%%%%%%%%%%%%%%%%%%%%%%%%%%%%%%%%%%%%%%%%%%%%%

%%%%%%%%%%%%%%%%%%%%%%%%%%%%%%%%%%%%%%%%%%%%%%%%%%%%%%%%%%%%%%%%%%%%%%%%
\begin{frame}{Model earnings elasticities}
\centering
\resizebox{0.7\textwidth}{!}{\input{tex/elasticities_by_type}}

\hyperlink{calibration}{\beamergotobutton{Back}}
\end{frame}
%%%%%%%%%%%%%%%%%%%%%%%%%%%%%%%%%%%%%%%%%%%%%%%%%%%%%%%%%%%%%%%%%%%%%%%%


%%%%%%%%%%%%%%%%%%%%%%%%%%%%%%%%%%%%%%%%%%%%%%%%%%%%%%%%%%%%%%%%%%%%%%%%
\begin{frame}{Internally calibrated parameters
}\label{table3}
\centering
{\input{tex/betascalibration}}

{\input{tex/mpcemp_table}}

\hyperlink{calibration}{\beamergotobutton{Back}}
\end{frame}
%%%%%%%%%%%%%%%%%%%%%%%%%%%%%%%%%%%%%%%%%%%%%%%%%%%%%%%%%%%%%%%%%%%%%%%%




%%%%%%%%%%%%%%%%%%%%%%%%%%%%%%%%%%%%%%%%%%%%%%%%%%%%%%%%%%%%%%%%%%%%%%%%
\begin{frame}{Heterogeneous effect on employment}\label{semp}
\begin{columns}[c]
\begin{column}{0.6\textwidth}
    \centering
\includegraphics[height = 0.7\textheight]{Figures/res_empl_log.pdf}
\end{column}
\hspace{-10mm}
\begin{column}{0.4\textwidth}
    \begin{itemize}
        \item Employment share rises most at the bottom
        \vspace{0.5cm}
        \item Top almost unaffected
        \vspace{0.5cm}
        
\hyperlink{decomposition}{\beamergotobutton{Back to decomposition}} 
    \end{itemize}
\end{column}
\end{columns}
\end{frame}
%%%%%%%%%%%%%%%%%%%%%%%%%%%%%%%%%%%%%%%%%%%%%%%%%%%%%%%%%%%%%%%%%%%%%%%%

%%%%%%%%%%%%%%%%%%%%%%%%%%%%%%%%%%%%%%%%%%%%%%%%%%%%%%%%%%%%%%%%%%%%%%%%
\begin{frame}{Probability of switching jobs}\label{switching}
\begin{columns}[c]
\begin{column}{0.6\textwidth}
    \centering
\includegraphics[height = 0.7\textheight]{Figures/trans_n_switch.pdf}
\end{column}
\hspace{-10mm}
\begin{column}{0.4\textwidth}
    \begin{itemize}
        \item Job switching probability rarely significantly affected
        \vspace{0.5cm}
        \item Slight effect at the top
        \vspace{0.5cm}
        
    \hyperlink{transprob}{\beamergotobutton{Back}}
    \end{itemize}
\end{column}
\end{columns}
\end{frame}
%%%%%%%%%%%%%%%%%%%%%%%%%%%%%%%%%%%%%%%%%%%%%%%%%%%%%%%%%%%%%%%%%%%%%%%%

%%%%%%%%%%%%%%%%%%%%%%%%%%%%%%%%%%%%%%%%%%%%%%%%%%%%%%%%%%%%%%%%%%%%%%%%
\begin{frame}{Incidence of monetary policy with unemployment benefits}\label{UI}
\begin{columns}[c]
\begin{column}{0.49\textwidth}
    \centering
    \alert{Earnings response with UI benefits} 
    
    \vspace{3pt}
    \includegraphics[width=\linewidth]{Figures/earn_n_totUI.pdf}
    \begin{itemize}
        \item Still downward sloping, less pronounced
    \end{itemize}
\end{column}
\begin{column}{0.49\textwidth}
    \centering
    \alert{Excluding extensive margin} 
    
    \vspace{3pt}
    \includegraphics[width=\linewidth]{Figures/earn_n_e2e_dashUI.pdf}
    \begin{itemize}
        \item Extensive margin still matters \hyperlink{decomposition}{\beamergotobutton{Back}}

    \end{itemize}
\end{column}
\end{columns}
\end{frame}
%%%%%%%%%%%%%%%%%%%%%%%%%%%%%%%%%%%%%%%%%%%%%%%%%%%%%%%%%%%%%%%%%%%%%%%%

%%%%%%%%%%%%%%%%%%%%%%%%%%%%%%%%%%%%%%%%%%%%%%%%%%%%%%%%%%%%%%%%%%%%%%%%
\begin{frame}{Long-run employment impact of monetary policy}\label{LRempl}
\vspace{-8mm}
\begin{align*}
    empl^q_{t+h} = \alpha_{x,h} + \gamma_{x,h} \Delta i_t + \theta_{x,h} X_t + \epsilon_{x,t}, \quad h = -6, \ldots, 36 
\end{align*}
\begin{columns}[c]
\begin{column}{0.49\textwidth}
    \centering
    \alert{Probability of employment} 
    
    \vspace{3pt}
    \includegraphics[width=\linewidth]{Figures/irf_u_employed_fa_n_1.pdf}
\end{column}
\begin{column}{0.49\textwidth}
    \centering
    \alert{Marginal effects by type} 
    
    \vspace{3pt}
    \includegraphics[width=\linewidth]{Figures/irf_u_employed_fa_n_3raw.pdf}
\end{column}
\end{columns}
\vspace{5pt}
\textbf{Higher long-run re-employment probability after expansionary policy}
\end{frame}
%%%%%%%%%%%%%%%%%%%%%%%%%%%%%%%%%%%%%%%%%%%%%%%%%%%%%%%%%%%%%%%%%%%%%%%%

%%%%%%%%%%%%%%%%%%%%%%%%%%%%%%%%%%%%%%%%%%%%%%%%%%%%%%%%%%%%%%%%%%%%%%%%
\begin{frame}{Long-run earnings (incl. UI) impact of monetary policy}\label{LRearn}
\vspace{-8mm}
\begin{align*}
    earn^q_{t+h} = \alpha_{x,h} + \gamma_{x,h} \Delta i_t + \theta_{x,h} X_t + \epsilon_{x,t}, \quad h = -6, \ldots, 36 
\end{align*}
\begin{columns}[c]
\begin{column}{0.49\textwidth}
    \centering
    \alert{Earnings including UI} 
    
    \vspace{3pt}
    \includegraphics[width=\linewidth]{Figures/irf_u_earnUI_n_1.pdf}
\end{column}
\begin{column}{0.49\textwidth}
    \centering
    \alert{Marginal effects by type} 
    
    \vspace{3pt}
    \includegraphics[width=\linewidth]{Figures/irf_u_earnUI_n_3raw.pdf}
\end{column}
\end{columns}
\vspace{5pt}
\textbf{Expansionary policy leads to higher long-run earnings after unemployment} \hyperlink{stock}{\beamergotobutton{Back}}
\end{frame}
%%%%%%%%%%%%%%%%%%%%%%%%%%%%%%%%%%%%%%%%%%%%%%%%%%%%%%%%%%%%%%%%%%%%%%%%


%%%%%%%%%%%%%%%%%%%%%%%%%%%%%%%%%%%%%%%%%%%%%%%%%%%%%%%%%%%%%%%%%%%%%%%%
\begin{frame}{Gini coefficient after expansionary monetary policy}\label{gini}
\begin{columns}[c]
\begin{column}{0.49\textwidth}
    \centering
    \alert{} 
    
    \vspace{3pt}
    \includegraphics[width=\linewidth]{Figures/gini_n_.pdf}
\end{column}
\begin{column}{0.49\textwidth}
    \centering
    \alert{Marginal effects by type} 
    
    \vspace{3pt}
    \includegraphics[width=\linewidth]{Figures/gini_n_UI.pdf}
\end{column}
\end{columns}
\vspace{5pt}
\textbf{Expansionary policy leads to lower inequality} \hyperlink{stock}{\beamergotobutton{Back}}
\end{frame}
%%%%%%%%%%%%%%%%%%%%%%%%%%%%%%%%%%%%%%%%%%%%%%%%%%%%%%%%%%%%%%%%%%%%%%%%

%%%%%%%%%%%%%%%%%%%%%%%%%%%%%%%%%%%%%%%%%%%%%%%%%%%%%%%%%%%%%%%%%%%%%%%%
\begin{frame}{More baseline model impulse responses}\label{baseirf1}
\begin{columns}[c]
\begin{column}{0.49\textwidth}
    \centering
    \alert{Unemployment} 
    
    \vspace{3pt}
    \includegraphics[trim=0cm 6cm 0cm 6.5cm, clip, width=\linewidth]{modelresults/single_Uagg_T_bgy397_tp0_ui6_sr0.pdf}
    \begin{itemize}
        \item Unemployment rises
    \end{itemize}
\end{column}
\begin{column}{0.49\textwidth}
    \centering
    \alert{Excluding extensive margin} 
    
    \vspace{3pt}
    \includegraphics[trim=0cm 6cm 0cm 6.5cm, clip, width=\linewidth]{modelresults/single_wTlev_bgy397_tp0_ui6_sr0.pdf}
    \begin{itemize}
        \item Real wage falls, forcing transitions \hyperlink{Cirf}{\beamergotobutton{Back}}
    \end{itemize}
\end{column}
\end{columns}
\end{frame}
%%%%%%%%%%%%%%%%%%%%%%%%%%%%%%%%%%%%%%%%%%%%%%%%%%%%%%%%%%%%%%%%%%%%%%%%


%%%%%%%%%%%%%%%%%%%%%%%%%%%%%%%%%%%%%%%%%%%%%%%%%%%%%%%%%%%%%%%%%%%%%%%%
\begin{frame}{Heterogeneous vs homogeneous incidence decomposition}\label{hetinc}
\begin{columns}[c]
\begin{column}{0.6\textwidth}
    \centering
    \includegraphics[trim=0cm 6cm 0cm 6.5cm, clip, width=\linewidth]{modelresults/aggs3_2_CTlev_bgy397_tp0_ui6_sr.pdf}
\end{column}
\hspace{-10mm}
\begin{column}{0.4\textwidth}
\begin{itemize}
    \item Keep Het U steady state \& apply homogeneous unemployment movements
        \vspace{0.3cm}

    \item Same labor input conditional on $w$
    
    \vspace{0.3cm}
    \item[$\to$] Heterogeneous transitions explain about 30\% of difference 

    \hyperlink{hetinc0}{\beamergotobutton{Back}}
    
\end{itemize}
\end{column}
\end{columns}


\end{frame}
%%%%%%%%%%%%%%%%%%%%%%%%%%%%%%%%%%%%%%%%%%%%%%%%%%%%%%%%%%%%%%%%%%%%%%%%


\end{document}
